% The interval scale itself suffices as the main Pell scale.
\section{Reusing the interval scale}
\label{sec:unit-scale}

There is another way to remove an operation from the base-four
construction. The scale $M=RY_P$ can be replaced by $M=Y_P$.
This saves its multiplication by $R$. The larger approximation error
is still small enough, provided it is bounded only after the
base-four exponent relation has been proved. A previously unused
upper bound on the packed index supplies the earlier estimates.

Start from Section~\ref{sec:main-base-four}, with the same fixed
modular-Sidon encoding and $L=2^{32}$. Set
\[
\begin{gathered}
 N=n,\quad R=r,\quad U=wN^2,\quad Y_P=sN^2,\quad M=Y_P,\\
 a_0=Y_P(U+1),\quad A=a_0+4,\quad
 Q=UY_P^2,\quad P=2Q+1,\quad J=2R+1.
\end{gathered}
\]
The supplied positive unknown $a$ continues to mean $a_0$.
Retain the first norm, affine index equation and positive interval:
\begin{equation}\label{eq:unit-scale-system}
\begin{gathered}
 \tau(\tau+1)=Q(Q+1)k^2,\qquad k=R+1+hQ,\\
 c=kY_P+\eta,\qquad k=\eta+\zeta,
 \qquad\tau,h,\eta,\zeta>0.
\end{gathered}
\end{equation}
Only the expanded source polynomials E9, E11 and E12 change, through
the definitions of $Q$ and $a_0$. The main parameter $A=a_0+4$,
signed auxiliary block and both exponent congruences retain their
forms in the supplied coordinates.

\subsection{A bound from the packing and the exact indices}

The pre-Pell coding proof gives
\[
 N=q^8,\quad q>4b\ge8,\quad 3L\le B\le N\le R,
 \quad b,q\le N,
 \quad0<S<N,\quad0\le T<N.
\]
The packed-index formula and integrality of the block values give
the useful additional bound
\begin{align}
 R&=S(N^2-N)+(T+1)(N^2-1)\notag\\
  &\le(N-1)(N^2-N)+N(N^2-1)
   =2N^3-2N^2<2N^3.\label{eq:unit-scale-rbound}
\end{align}
Thus $N,R\ge64$ and $U,Y_P\ge N^2\ge4096$. In particular
\begin{equation}\label{eq:unit-scale-initial}
 a_0>N^4>4N^3+1>J,\qquad
 Q\ge N^6>2N^3+1>R+1.
\end{equation}
These conclusions use only the coding bounds and positivity, before
either exponent relation is known.

The index equation gives $k\ge R+2$, and the interval gives
$c>Y_Pk>J$. The signed Pell block therefore implies
$c=\psi_A(J)$ and $d=\chi_A(J)$.
The first norm gives the ordinary Pell equation with root
$\widehat\tau=2\tau+1$ and parameter $P=2Q+1$, so
$k=\psi_P(t)$ for a positive index $t$.
The Pell congruence modulo $P-1=2Q$, reduced modulo $Q$, and
\eqref{eq:unit-scale-initial} give
\[
 t=R+1+vQ,\qquad v\ge0.
\]
If $v\ge1$, then Pell growth yields
\[
 \frac ck\le\frac{(2A)^{2R}}{(2P-1)^{R+Q}}
 \le\left(\frac{2A}{2P-1}\right)^{2R}<\frac12.
\]
Here $Q\ge R$, and the last inequality follows from
\[
 (2P-1)-4A=4Y_P\bigl(U(Y_P-1)-1\bigr)-15>0.
\]
This contradicts $c/k>Y_P$. Consequently
\begin{equation}\label{eq:unit-scale-indices}
 c=\psi_A(J),\quad d=\chi_A(J),\quad
 k=\psi_P(R+1),\quad\widehat\tau=\chi_P(R+1).
\end{equation}

\subsection{The larger error is bounded after exponent recovery}

Put $\xi=(U+1)^{2R}/U^R$. At the exact indices, the same cancellation
of $M$ in the principal ratio gives
\begin{equation}\label{eq:unit-scale-lower}
 \frac ck\ge
 \xi\left(1+\frac7{2a_0}\right)^{2R}
     \left(1+\frac1{2Q}\right)^{-R}>\xi,
\end{equation}
because $14Q>a_0$. The upper growth bound gives
\[
 \frac ck<\xi\left(1+\frac4{a_0}\right)^{2R}.
\]
The packed-index bound replaces the old factor $R$ in the scale:
\[
 \frac{8R}{a_0}<\frac{16N^3}{N^4}
 =\frac{16}{N}\le\frac14<\frac12.
\]
The elementary estimate $(1+t)^m\le(1-mt)^{-1}<1+2mt$
therefore gives
\begin{equation}\label{eq:unit-scale-upper}
 \frac ck<\xi\left(1+\frac{16R}{a_0}\right).
\end{equation}
The interval and the strict lower bound give $\xi<Y_P+1$.
Since $Y_P$ is integral and $\xi>U^R$, this implies
\begin{equation}\label{eq:unit-scale-size}
 Y_P\ge U^R,\qquad A>a_0=Y_P(U+1)>U^{R+1}.
\end{equation}
Also $\xi<Y_P+1<2Y_P$, so
\begin{equation}\label{eq:unit-scale-error-unbounded}
 0<\frac ck-\xi<\frac{32R}{U+1}.
\end{equation}
At this point the right-hand side is not assumed to be less than one.

Write $W=bw$. Since $W\le U$, $N\ge64$ and $U\ge4096$,
\eqref{eq:unit-scale-size} gives
\[
 4^{3J}\le64N^{2R}\le64U^R<U^{R+1}<A,
 \qquad W^3\le U^3\le U^R<A.
\]
The base-four exponent congruence now proves $W=4^J$ by source
Lemma~2.22. Hence $b$ is a power of two and
\[
 U=N^2w\ge Nbw=N4^J\ge64\cdot4^J>64R.
\]
Only now does~\eqref{eq:unit-scale-error-unbounded} yield
\begin{equation}\label{eq:unit-scale-small-error}
 0<\frac ck-\xi<\frac12.
\end{equation}

The gap, second Pell norm and index congruence give
$\kappa=\psi_A(L)$ and $\mu=\chi_A(L)$, since
$0<L<J<A$. Furthermore
\[
 B^{3L}\le B^B\le N^N\le U^R<A,
 \qquad q^3\le N^N\le U^R<A.
\]
The second exponent congruence therefore gives $q=B^L$.
Since also $U\ge N4^J>4\cdot2^{2R}$, the binomial expansion has
integer part $F=\lfloor\xi\rfloor$ with
\[
 0<\xi-F<\frac14,\qquad
 F\equiv\binom{2R}{R}\pmod U.
\]
The two inequalities
$F<\xi<c/k<Y_P+1$ and
$Y_P<c/k<\xi+1/2<F+3/4$
force $Y_P=F$. Thus $N^2\mid\binom{2R}{R}$, and the unchanged
no-carry and modular-Sidon decoding proves membership.

\subsection{Positive witnesses and the independent count}

For necessity, take the canonical coding witnesses, which determine
$b,B,q,N,R$ and satisfy~\eqref{eq:unit-scale-rbound}. Choose
\[
 J=2R+1,\quad w=4^J/b,\quad U=N^2w,\quad
 Y_P=\lfloor\xi\rfloor,\quad s=Y_P/N^2,
\]
and then form $a_0=Y_P(U+1)$, $A=a_0+4$, $Q=UY_P^2$,
$P=2Q+1$, $c=\psi_A(J)$, $d=\chi_A(J)$ and
$k=\psi_P(R+1)$.
The positive integrality of $w,s$ follows from the same power-of-two
and binomial arguments as in Section~\ref{sec:main-base-four}.
The bounds~\eqref{eq:unit-scale-lower}--\eqref{eq:unit-scale-upper}
apply at these indices, and $U\ge N4^J$ holds directly from the
chosen $w$. Hence
\[
 Y_P<\xi<\frac ck<\xi+\frac12<Y_P+\frac34.
\]
Choose the positive integers
\[
 \eta=c-kY_P,\qquad \zeta=k-\eta,\qquad
 \tau=\frac{\chi_P(R+1)-1}{2},\qquad
 h=\frac{k-R-1}{Q}.
\]
The oddness of $P$ gives the integrality of $\tau$, and the Pell
congruence modulo $2Q$ gives that of $h$. Strict Pell growth gives
their positivity.

All remaining dependent Pell witnesses are obtained by the generic
constructions of Section~\ref{sec:main-base-four} at this new $A$.
Explicitly, choose positive $f,i$ for E16 and put
\[
 G=1+(A+1)(f^2-1),\qquad
 o=\frac{\chi_G(J)+d}{f},\qquad
 j=\frac{\psi_G(J)-J}{c}.
\]
The congruences $G\equiv-A\pmod f$, $G\equiv1\pmod c$ and oddness
of $J$ make $o,j$ positive integers. Choose
$\kappa=\psi_A(L)$, $\mu=\chi_A(L)$,
$\Delta=(\kappa-L)/(A-1)$ and $\phi=c-\kappa$; these are
positive integers because $2\le L<J$.
The exponent quotients are positive since
\[
 d-(A-4)c>3c>4^J,\qquad
 \mu-(A-B)\kappa>(B-1)\kappa>B^L,
\]
using $A>(4^J)^3,(B^L)^3$ and $J,L\ge2$.
Their moduli are positive. All coding witnesses remain unchanged.
This constructs every required positive coordinate at the new scale.

\begin{theorem}\label{thm:unit-scale105}
For each r.e.\ set, reusing $Y_P$ as the Pell scale gives a
membership-equivalent system in 34 positive unknowns and 22 equations
with a straight-line certificate of \textbf{105 arithmetic instructions}:
58 multiplications and 47 additions. Fixed numerals, supplied parameters
and equality tests are free.
\end{theorem}
\begin{proof}
The preceding arguments establish soundness and positive necessity.
The predecessor already computes $Y_P=sN^2$. Delete the multiplication
$M=RY_P$ and reuse $Y_P$ wherever $M$ was used. The three instructions
\[
 UY_P=U\cdot Y_P,\qquad a_0=UY_P+Y_P,
 \qquad Q=(UY_P)Y_P
\]
retain their previous costs. The affine index equation retains its
multiplication and addition. Thus exactly one multiplication disappears
from the 106-operation certificate. The complete residual and primitive
checks are supplied by
\file{round18\_1980\_unit\_scale\_certificate.py}.
\end{proof}

This is independent of Theorem~\ref{thm:best105}: its saving is a
multiplication rather than an addition. The index-multiple argument
there cannot be applied automatically here, since $P-1=2UY_P^2$ is
not known to be divisible by $R+1$.

\subsection{Sharing the interval product with the norm}

The smaller scale permits another multiplication to be removed.
Let
\[
 D=UY_P,\qquad Q=DY_P=UY_P^2,\qquad H_P=Y_Pk.
\]
Both $D$ and $H_P$ are already computed: $D$ is needed for
$a_0=D+Y_P$, and $H_P$ is needed for $c=H_P+\eta$.
The polynomial identity
\begin{equation}\label{eq:shared-pell-norm}
 Q(Q+1)k^2
 =(UY_P^2)(UY_P^2+1)k^2
 =\bigl((UY_P)^2+U\bigr)(Y_Pk)^2
 =(D^2+U)H_P^2
\end{equation}
allows the norm to share these two products.
To remove the remaining use of the computed value $Q$, replace
the index equation in~\eqref{eq:unit-scale-system} by
\begin{equation}\label{eq:shared-pell-index}
 k=R+1+hD,\qquad h>0.
\end{equation}
The norm and every other source polynomial are unchanged.

We first prove that the smaller modulus still gives the exact index.
Before exponent decoding, the bounds above give
\[
 D=UY_P\ge N^4>2N^3+1>R+1,\qquad A>J.
\]
The new equation gives $k\ge R+2$, so the interval gives $c>J$.
The signed Pell block gives $c=\psi_A(J)$ and $d=\chi_A(J)$.
The unchanged triangular norm gives $k=\psi_P(t)$ for some
positive index $t$, where $P=2Q+1$. Since
$P-1=2Q=2DY_P$ is divisible by $D$, the Pell congruence and
\eqref{eq:shared-pell-index} imply
\[
 t=R+1+vD,\qquad v\ge0.
\]
If $v\ge1$, then $D>R$ and the same growth estimate gives
\[
 \frac ck\le\frac{(2A)^{2R}}{(2P-1)^{R+D}}
 \le\left(\frac{2A}{2P-1}\right)^{2R}<\frac12,
\]
contradicting $c/k>Y_P$. Therefore $t=R+1$.
This recovery uses only the preliminary coding bounds, the interval
and the two Pell blocks.

There are direct maps between the positive witnesses of this system
and those of Theorem~\ref{thm:unit-scale105}. From an old solution,
put $h=Y_Ph_{\rm old}$ and preserve every other witness. Then
$hD=h_{\rm old}Q$, so the new index equation holds.
Conversely, the recovered exact index and the stronger Pell congruence
modulo $2Q$ show that
\[
 h_{\rm old}=\frac{k-R-1}{Q}
\]
is an integer. It is positive because $\psi_P(R+1)>R+1$.
It restores the old index equation, preserving every other witness.
Equivalently, the new $h$ is a positive multiple of $Y_P$ and the
reverse map is $h_{\rm old}=h/Y_P$.
This proves equivalence over the positive domains before invoking
the unit-scale approximation, exponent or rounding conclusions.
Those conclusions and the complete decoding proof therefore transfer.

\begin{theorem}\label{thm:best104}
For each r.e.\ set, the fixed modular-Sidon encoding gives a
membership-equivalent system in 34 positive unknowns and 22 equations
whose straight-line certificate has \textbf{104 arithmetic instructions}:
57 multiplications and 47 additions. Fixed numerals, supplied parameters
and equality tests are free.
\end{theorem}
\begin{proof}
Positive-domain equivalence follows from the preceding witness maps
and Theorem~\ref{thm:unit-scale105}. For the count, move the existing
computation of $H_P=Y_Pk$ before the norm. The former five operations
\[
 Q=DY_P,\quad Q_+=Q+1,\quad F_Q=QQ_+,
 \quad k_2=k^2,\quad L_9=F_Qk_2
\]
are replaced by four:
\[
 D_2=D^2,\quad F_D=D_2+U,\quad H_2=H_P^2,
 \quad L_9=F_DH_2.
\]
Identity~\eqref{eq:shared-pell-norm} proves that they give the same
norm residual. The computation of $H_P$ is moved rather than copied,
and the root-side instructions $\tau+1$ and $\tau(\tau+1)$ are
unchanged. Both displayed lists contain one addition; the second
contains three multiplications rather than four.

The product $h_{\rm old}Q$ is replaced by $hD$ at the same cost.
Thus the calculated register $Q$ is no longer required anywhere in
the certificate. Its mathematical definition in the proof introduces
no additional supplied witness or uncharged computed register.
The total is 57 multiplications and 47 additions. The checker
\file{round19\_1980\_shared\_ratio\_product\_certificate.py}
verifies the identity, the forward witness-map polynomial, all source
residuals and primitive instructions, and the absence of the discarded
registers. It checks explicitly that only E11 changes as a source
polynomial.
\end{proof}
